% 第5章 5.1.2：三种核在 [-1,1] 上的核截面 K_x(t)=K(t,x)，x=-0.6,0,0.6（解析式）。
% (a) 高斯核 exp(-(t-x)^2/(2*0.3^2))；(b) 仿射核 1+xt；(c) 二次多项式核 (1+xt)^2。
\begin{tikzpicture}[font=\figurefont\footnotesize]
\pgfplotsset{every axis/.style={width=40mm,height=34mm,axis lines=left,xmin=-1,xmax=1,
  xtick={-1,0,1},
  x axis line style={draw=FigureStroke,line width=.6pt},y axis line style={draw=FigureStroke,line width=.6pt},
  tick style={draw=FigureStroke},tick label style={font=\figurefont\scriptsize,text=FigureInk},
  label style={font=\figurefont\scriptsize},xlabel={$t$},xlabel style={at={(axis description cs:1,0)},anchor=north west,yshift=2mm},
  ymajorgrids,grid style={draw=FigureGrid,line width=.4pt},
  title style={font=\figurefont\footnotesize},samples=121,domain=-1:1}}
\begin{axis}[name=A,ymin=0,ymax=1.1,ytick={0,0.5,1},title={(a) 高斯核}]
\addplot[draw=FigureBlue,line width=1.1pt] {exp(-(x+0.6)^2/0.18)};
\addplot[draw=FigureGreen,line width=1.1pt,dashed] {exp(-(x)^2/0.18)};
\addplot[draw=FigureYellow,line width=1.1pt,dash dot] {exp(-(x-0.6)^2/0.18)};
\end{axis}
\begin{axis}[name=B,at={(A.east)},anchor=west,xshift=8mm,ymin=0,ymax=1.7,ytick={0,0.5,1,1.5},title={(b) 仿射核},
  legend style={draw=none,fill=none,font=\figurefont\scriptsize,at={(0.5,-0.32)},anchor=north,legend columns=3,/tikz/every even column/.append style={column sep=3mm}}]
\addplot[draw=FigureBlue,line width=1.1pt] {1-0.6*x};
\addplot[draw=FigureGreen,line width=1.1pt,dashed] {1};
\addplot[draw=FigureYellow,line width=1.1pt,dash dot] {1+0.6*x};
\legend{$x=-0.6$,$x=0$,$x=0.6$}
\end{axis}
\begin{axis}[at={(B.east)},anchor=west,xshift=8mm,ymin=0,ymax=2.7,ytick={0,1,2},title={(c) 二次多项式核}]
\addplot[draw=FigureBlue,line width=1.1pt] {(1-0.6*x)^2};
\addplot[draw=FigureGreen,line width=1.1pt,dashed] {1};
\addplot[draw=FigureYellow,line width=1.1pt,dash dot] {(1+0.6*x)^2};
\end{axis}
\end{tikzpicture}
